\begin{figure}[tbp]
\centering
\begin{tikzpicture}[
  box/.style={draw=blue!45!black,rounded corners=2pt,fill=blue!3,
    text width=5.6cm,minimum height=1.2cm,align=center,
    inner sep=5pt,font=\small},
  flow/.style={-{Stealth[length=2mm]},thick,draw=blue!45!black}
]
\node[box] (length) {Fix the number of colors $r$.\\
  Choose a Hales--Jewett length $n$.};
\node[box,right=9mm of length] (endpoint) {Prepare one endpoint $w$\\
  with paths of weights $1,\frac12,\ldots,\frac1n$.};
\node[box,below=10mm of endpoint] (line) {Color words using $w$.\\
  A monochromatic line has $t$ variable positions.};
\node[box,left=9mm of line] (scale) {Use the path of weight $1/t$\\
  at each variable position.\\
  Squared-distance factor: $t(1/t)=1$.};
\draw[flow] (length) -- (endpoint);
\draw[flow] (endpoint) -- (line);
\draw[flow] (line) -- (scale);
\end{tikzpicture}
\caption{The order of choices in the path criterion. Every possible
variable count is accommodated before the word coloring is defined.
Choosing a path representation later leaves its endpoint unchanged;
common factors become common strings, so only the variable blocks
contribute to pairwise distances.}
\label{fig:synchronization}
\end{figure}
